% 05_conclusion.tex — Conclusion
% Drafted Session 4bd — block-and-wait review gate.

\section{Conclusion}\label{sec:conclusion}

We have shown that a cheap, reference-anchored global-optimization screening
identifies amorphous minima in an interstitial alloy. The method models the
amorphous state as the low-energy basin around a chosen crystalline reference,
enumerates that basin with a surrogate-guided search, and confirms
amorphization with the Steinhardt $q_6$ crystalline fraction; the reference is
a physical choice the user controls, unique for a host with a single stable
phase and a genuine choice for a polymorphic host. Applied to interstitial
Pt(P), the screening finds that 20--30~at.\% P marks the onset of structural
amorphization: the crystalline fraction collapses from 0.99 in the undoped host
to 0.09--0.18, a trend cross-checked by the XRD integrated crystallinity index.
The effect saturates with concentration (the 20P and 30P systems are
indistinguishable) and is mildly promoted by cell expansion, which adds free
volume and acts synergistically with the P-induced frustration of fcc packing,
all at a fraction of the cost of first-principles melt-quench molecular
dynamics.
